\documentclass[border=5pt]{standalone}
\usepackage{pgfplots}
\usepackage{pgf-pie} % 饼图：AI 代码可直接使用 \pie
\pgfplotsset{compat=1.18}
\usepgfplotslibrary{fillbetween}  % 面积图 / 堆叠面积图 / \closedcycle 填充路径
\usepgfplotslibrary{errorbars}     % \addplot +[error bars/.cd, ...] 误差棒子键
\usepackage{amsmath}
\usepackage{amssymb}
\usepackage{xcolor}

% 支持中文
\usepackage{fontspec}
\usepackage{xeCJK}
\setCJKmainfont{SimSun}
\setmainfont{Times New Roman}

\begin{document}

\begin{tikzpicture}
\begin{axis}[
    xbar,
    width=11cm, height=8.5cm,
    bar width=0.5,
    title={2023年末主要城市轨道交通运营里程对比},
    xlabel={运营里程（km）},
    ylabel={城市},
    symbolic y coords={南京,武汉,杭州,深圳,广州,成都,北京,上海},
    ytick=data,
    xmin=0, xmax=950,
    enlarge y limits=0.08,
    nodes near coords,
    every node near coord/.append style={font=\scriptsize, fill=none, draw=none, inner sep=1pt, anchor=west},
    legend style={at={(1.03,0.5)}, anchor=west, draw=black, fill=white}
]
\addplot[fill=blue!45, draw=blue!70!black] coordinates {
    (831,上海)
    (807,北京)
    (632,成都)
    (621,广州)
    (567,深圳)
    (516,杭州)
    (486,武汉)
    (459,南京)
};
\legend{轨道交通运营里程}
\node[anchor=north west, font=\scriptsize] at (axis description cs:0.0,-0.15) {数据来源：2023年城市轨道交通运营数据（交通运输部）};
\end{axis}
\end{tikzpicture}

\end{document}